% Job posting: https://jobs.ashbyhq.com/eli/dd8575a8-60e0-472f-a8b9-6f7a6478699d
% =====================================================================
%  CV - Nishal Stanislaus Silva, PhD
%  Target: Eli Health - Machine Learning Scientist, Montreal QC (on-site)
%  Variant: V1 (AI-Industry) with an applied-science / statistical-rigor lean
% =====================================================================

\documentclass[11pt]{article}
\usepackage[left=0.7in,top=0.7in,right=0.7in,bottom=0.7in]{geometry}
\usepackage{enumitem}
\usepackage[hidelinks]{hyperref}
\usepackage{titlesec}
\usepackage{array}
\usepackage{setspace}
\usepackage{needspace}
\usepackage[T1]{fontenc}
\usepackage{newtxtext,newtxmath}
\ifdefined\pdfgentounicode
  \input{glyphtounicode}
  \pdfgentounicode=1
\fi
\setstretch{1.05}
\pagenumbering{gobble}
\titleformat{\section}{\Large\scshape}{}{4em}{}[\vspace{-0.7em}\rule{\linewidth}{0.4pt}\vspace{-0.4em}]
\titlespacing*{\section}{0pt}{1em}{0.4em}
\newenvironment{cvrole}[5]{%
  \par\noindent\textbf{\large #1} | \textit{\small #2, #3}\hfill \textit{\small #4 - #5}\par
  \begin{itemize}[leftmargin=2em, itemsep=-0.3em, topsep=-0.5em]%
}{%
  \end{itemize}\par\noindent\mbox{}\par\vspace{0.1em}%
}
\newcommand{\cvName}{Nishal Stanislaus Silva}
\newcommand{\cvEmail}{nishal.silva@hotmail.com}
\newcommand{\cvPhone}{+1 613 200 2030}
\newcommand{\cvCity}{Toronto, ON}
\newcommand{\cvSite}{https://nishal.xyz}
\newcommand{\cvLinkedIn}{https://www.linkedin.com/in/nishal-silva}
\newcommand{\cvGitHub}{https://github.com/ns2max}
\newcommand{\cvStatus}{Canadian Permanent Resident, Eligible to work in Canada without sponsorship}
\newcommand{\cvTagline}{Machine Learning Research Engineer | Deep Learning, Time-Series, Pattern Detection | Applied ML}

\begin{document}
\pagestyle{empty}
\begin{center}
    {\LARGE \scshape \textbf{\cvName}}{\large \textit{, PhD}}\\[1pt]
    {\small \cvTagline}\\[1pt]
    {\footnotesize\textit{\cvStatus}}\\[1pt]
    {\small
    \href{mailto:\cvEmail}{\cvEmail} \,|\, \cvPhone \,|\, \cvCity
    \,|\, \href{\cvSite}{nishal.xyz} \,|\, \href{\cvLinkedIn}{linkedin.com/in/nishal-silva}
    \,|\, \href{\cvGitHub}{github.com/ns2max}}
\end{center}

\section*{Professional Summary}

Machine learning research engineer (PhD, University of Trento, 2025) and 10+ years across industry and academia, experienced on turning imperfect, noisy real-world sensor data into validated, reproducible models. My work centers on signal processing and statistical modeling on messy data, model validation/generalization, and benchmark and ablation design. During my postdoc I designed the frozen-protocol benchmarks behind a published detector running at F1 0.76 in 14 ms on a Raspberry Pi 4, 74$\times$ faster than the DTW baseline it replaced. My earlier work carried the same discipline into deployed systems: financial anomaly detection on messy structured data at Forestpin, and computer-vision inspection shipped across factories and COVID-19 remote-vitals detection system for hospital wards. Experienced in diagnosing production failures and distribution shift from real-world data and communicating findings to technical and non-technical audiences. 

\section*{Core Competencies}

\textbf{ML Engineering:} signal processing and feature engineering on noisy data (DSP, STFT, MFCC, S-transform), statistical and probabilistic modelling, model selection and validation, bias/variance and generalization analysis, uncertainty quantification, classical ML and deep sequence models, anomaly detection, reproducible Python beyond notebooks (pytest, CI/CD)\\
\textbf{Domain Expertise:} streaming time-series sensor data (audio, IMU, EEG/BCI, XR), multimodal sensor fusion and multi-stream synchronization, medical-signal adjacency (COVID-19 remote-vitals CV device across hospital wards), financial anomaly detection, industrial computer vision and IoT at scale\\
\textbf{Research \& Leadership:} benchmark and ablation design, evaluation under domain shift, error analysis and production-failure diagnosis, experiment design to resolve uncertainty, communicating findings to technical and non-technical audiences, mentoring and group-wide training

\section*{Technical Stack}

\textbf{Languages:} Python, C++, C, MATLAB, SQL, Shell\\
\textbf{ML \& Statistics:} TensorFlow, Keras, scikit-learn, PyTorch, NumPy, Pandas, SciPy; RNN/LSTM/GRU, CNN, anomaly detection, Youden-index thresholding, ablation and uncertainty analysis, VAE and diffusion prototypes\\
\textbf{Signal \& Data:} DSP, STFT, MFCC, S-transform, Librosa, SciPy, FFTW; OpenCV, image registration, feature extraction; ETL pipelines, dataset and benchmark design\\
\textbf{Dev \& Infra:} Docker, Airflow, Jenkins, Git, pytest, CI/CD, REST APIs, AWS (EC2, S3, ECS, MWAA, RDS)

% \section*{Education}
% \textbf{Ph.D - Information Engineering and Computer Science} \textit{\small{ - University of Trento, Italy}} \hfill \textit{Jan 2025}\\[2pt]
% \noindent\textit{\small Thesis: Embedded Real-time Musical Pattern Detection for Smart Musical Instruments}\\[3pt]
% \noindent\textbf{M.Sc - Telecommunication and Electronic Engineering} \textit{\small{ - Sheffield Hallam University, UK}} \hfill \textit{Aug 2018}\\[3pt]
% \noindent\textbf{B.Eng (Hons) - Electronic Engineering} \textit{\small{ - Sheffield Hallam University, UK}} \hfill \textit{Mar 2014}

\section*{Selected Projects \footnotesize {\normalfont{(full portfolio: \href{https://nishal.xyz/research}{\textit{nishal.xyz/research}})}}}
\begin{itemize}[leftmargin=1.2em, itemsep=-0.25em]
    \item \textbf{Real-time audio pattern detection:} MFCC front end into a stacked RNN on 30 ms windows; Raspberry Pi 4 at 14 ms, 31.4\% CPU, F1 0.76 against a DTW baseline at 1038 ms.
    \item \textbf{SSIM-based training-free detection:} four-attribute representation with bicubic resizing for variable length; 95\% detection across human, synthetic and JKUPDD sets with no training data.
    \item \textbf{Open benchmark datasets:} four open-access musical-pattern datasets, 9,800+ recordings from 70+ musicians, published on Zenodo with frozen evaluation protocols.
\end{itemize}


\newpage

\section*{Professional Experience}

\begin{cvrole}{Independent ML Researcher}{Self-directed}{Toronto, Canada}{Jan 2026}{Present}
    \item MIRaaS latency-characterization paper accepted at IEEE IS\textsuperscript{2} 2026; Smart Drums paper in press at JAES; Released Nebula, an open-source C++ feature extraction library.
\end{cvrole}

\needspace{5\baselineskip}
\begin{cvrole}{Postdoctoral Researcher}{University of Trento}{Trento, Italy}{Jan 2025}{Dec 2025}
    \item Owned full ML pipelines for streaming audio and IMU/gesture data: acquisition, DSP and feature engineering, training and evaluation, edge deployment with monitoring; EU Horizon EIC Pathfinder project MUSMET.
    \item Designed frozen-protocol benchmark suites and ablations (RNN vs DTW; audio vs MIDI; pattern-set sizes) to separate real accuracy gains from latency and CPU noise; published results F1 0.76 at 14 ms on a Raspberry Pi 4.
    \item Built a real-time multimodal synchronization pipeline aligning audio, BCI/EEG and mixed-reality head/hand tracking across multiple personnel to create a comprehensive, multidimentional dataset.  
    % musicians simultaneously, with SAE Institute Barcelona; ran a user study across two live concerts (20 audience, 6 performers, p $<$ 0.05).
\end{cvrole}

\needspace{5\baselineskip}
\begin{cvrole}{Visiting Researcher}{McGill University (IDMIL / CIRMMT)}{Montreal, Canada}{May 2024}{Aug 2024}
    \item Fused inertial (IMU) and audio streams through a stacked deep-learning and hierarchical finite-state machine to develop a multimodal pattern detection system; profiled compute and power for real-time feasibility.
    % accumulating repeated observations of the same event to reach F1 0.78 across a 500-event corpus (IEEE IS\textsuperscript{2} 2025); profiled compute and power for real-time feasibility.
\end{cvrole}

\needspace{5\baselineskip}
\begin{cvrole}{Technical Consultant}{Forestpin (Pvt) Ltd}{Colombo, Sri Lanka}{Aug 2020}{Feb 2021}
    \item Designed and deployed ML and statistical pipelines for financial forensics and risk detection on large, messy structured enterprise datasets; SQL and Python backend anomaly-scoring services in production.
\end{cvrole}

\needspace{5\baselineskip}
\begin{cvrole}{Research Engineer}{MAS Holdings (Pvt) Ltd}{Colombo, Sri Lanka}{Jan 2015}{Jul 2020}
    \item Built and deployed end-to-end computer-vision and IoT systems for manufacturing at scale, integrating camera, sensor and IoT streams into live production workflows; improved operational efficiency by up to 300\%.
    \item Shipped automated QC pipelines (OpenCV template and feature matching, conveyor with PLC reject); inspection time cut 99.5\% with a human-in-the-loop path for borderline cases.
    % \item Delivered the COVID-19 remote-vitals device: a computer-vision system digitizing patient vitals from retrofitted webcams, deployed across hospital wards for remote monitoring (GMOA recognition).
    \item Introduced CI/CD and code review, mentored engineers and ran group-wide computer-vision training; delivered an on-demand garment-printing system (geometry-alignment algorithm + industrial integration).
\end{cvrole}


\section*{Education}
\textbf{Ph.D - Information Engineering and Computer Science} \textit{\small{ - University of Trento, Italy}} \hfill \textit{Jan 2025}\\[2pt]
% \noindent\textit{\small Thesis: Embedded Real-time Musical Pattern Detection for Smart Musical Instruments}\\[3pt]
\noindent\textbf{M.Sc - Telecommunication and Electronic Engineering} \textit{\small{ - Sheffield Hallam University, UK}} \hfill \textit{Aug 2018}\\[2pt]
\noindent\textbf{B.Eng (Hons) - Electronic Engineering} \textit{\small{ - Sheffield Hallam University, UK}} \hfill \textit{Mar 2014}

\enlargethispage{3\baselineskip}
\section*{Publications \& Recognition}

\begin{itemize}[leftmargin=1.2em, itemsep=-0.25em]

    \item \textbf{GMOA Recognition (Sri Lanka):} delivered a computer-vision device digitizing patient vitals from retrofitted webcams, deployed across hospital wards for COVID-19 remote monitoring.
    \item 10+ peer-reviewed papers (JAES, IEEE, ACM) {\normalfont{Complete publications portfolio: \href{https://nishal.xyz/publications}{\textit{nishal.xyz/publications}})}}.
    \item Released four open-access Zenodo datasets; MIDI Innovation Awards 2023 finalist.
    \item Released Nebula, a C++17 feature-extraction library with per-feature ARM latency benchmarks.
    \item Currently researching foundation model based generation for low-training data regimes.
    
    % incl. \textit{Real-Time Audio Pattern Detection for Smart Musical Instruments}, J.~Audio Eng.~Soc.~74(3), 2026 (F1 0.76 at 14 ms on a Raspberry Pi 4) and the DAFx 2022 training-free method (95\% detection); four open-access Zenodo datasets (9,800+ recordings, 70+ musicians); MIDI Innovation Awards 2023 finalist. Full list: \href{https://nishal.xyz/publications}{nishal.xyz/publications}.


    % \item \textbf{Research and open tooling:} 10+ peer-reviewed papers (JAES, IEEE, ACM); Nebula, a C++17 feature-extraction library with per-feature ARM latency benchmarks; four open-access Zenodo datasets; {\normalfont{Complete publications portfolio: \href{https://nishal.xyz/publications}{\textit{nishal.xyz/publications}})}}
\end{itemize}

\end{document}
